The remaining zero-gap case has one supercritical V triangle. There is no
boundary gap to use as a witness, so the obstruction must be built inside
$H$. The construction and its verification have three stages.

\readerheading{Force the points.}
Strict boundary handoffs give common lower bounds for the boundary reaches.
These force six radial points into the C triangle. Three further points are
chosen on the exact frontier of the feasible union of the supercritical
triangle; separate diameter and handoff arguments exclude every other V
triangle. Exclusion from the supercritical triangle alone would not be enough.
A frontier point may belong to a closed feasible triangle while being missed
by its open interior.

\readerheading{Simplify the set to be enclosed.}
The convex hull of the six radial points contains a disk. Replacing those
six points by this disk gives a smaller comparison set. In the nontrivial
range, the two outer frontier points are then moved inward by a single
Newton step at the fixed line parameter $k=1/2$. Convexity keeps these comparison points inside any candidate
containing the original witnesses. These are exact inner comparisons,
not numerical approximations to a covering configuration.

\readerheading{Bound the enclosing triangle.}
The inner comparison reduces the enclosure problem to four supporting
contacts: two contacts along consecutive-point lines and two disk-tangent
contacts. The line estimates are analytic. The remaining paired tangent
inequalities are reduced to the exact polynomial certificate. At the end,
the comparison set requires side at least one, contradicting its containment
in the open C triangle. The calculations below justify each of these stages;
the certificate enters only at the last one.

\readerheading{Active tangent calculation.}
The two outer comparison points use the fixed-start formulas, not the older
junction-start formulas. The exact coefficient comparisons and radius envelope
are specified in \zcref{thm:paper-exact-mixed-certificate}.
The actual disk is unchanged; the smaller scalar radius is used only for
initially checking the tangent residuals.
